% Experiment & Result Section
% Earthquake Prediction Problems - Data Science Report

This section presents the experimental setup, model implementations, and results for three machine learning tasks: magnitude regression, depth category classification, and seismic hotspot clustering.

\subsection{Experimental Setup}

\subsubsection{Computing Environment}

All experiments were conducted using Python 3.10 with standard data science libraries. The primary libraries used include Pandas 2.0 for data manipulation, NumPy 1.24 for numerical computing, Scikit-learn 1.3 for machine learning algorithms, and Matplotlib 3.7 with Seaborn 0.12 for visualization. Experiments were executed on a system with 16GB RAM and an Intel Core i7 processor.

\subsubsection{Data Preparation}

The dataset was prepared following the feature selection process described in Section 6. The final dataset contains 2,886,684 records after removing instances with missing values in critical features. The data was split into training (70\%), validation (15\%), and test (15\%) sets using stratified sampling to maintain class proportions for classification tasks.

Due to the large dataset size, a stratified sample of 500,000 records was used for initial model development and hyperparameter tuning. Final models were trained on the full training set to maximize performance. All numerical features were standardized using z-score normalization with parameters computed from the training set only to prevent data leakage.

\subsection{Task 1: Magnitude Regression}

\subsubsection{Task Definition}

The objective of this task is to predict the exact magnitude value of earthquakes using regression techniques. Unlike categorical classification, regression provides continuous magnitude estimates that enable finer-grained risk assessment and early warning applications.

\subsubsection{Input Features}

The input features for magnitude regression include latitude representing the north-south epicenter position, longitude representing the east-west epicenter position, depth of the earthquake hypocenter in kilometers, azimuthal gap (gap) indicating the quality of station coverage, and number of recording stations (nst). These features were selected based on their correlation with magnitude observed during exploratory data analysis.

\subsubsection{Model Implementation}

Random Forest Regressor was implemented as the primary algorithm for this task. Random Forest is an ensemble method that combines multiple decision trees through bootstrap aggregation (bagging) to reduce variance and improve generalization. The model naturally handles nonlinear relationships and feature interactions without requiring extensive feature engineering.

The model was configured with 100 estimators (trees) in the ensemble, maximum tree depth of 20 to prevent overfitting, minimum samples per leaf of 10 for regularization, and the square root of total features considered at each split. Hyperparameters were tuned using 5-fold cross-validation on the training set.

\subsubsection{Evaluation Metrics}

Two primary metrics were used to evaluate regression performance. Root Mean Square Error (RMSE) measures the average prediction error in magnitude units, calculated as $RMSE = \sqrt{\frac{1}{n}\sum_{i=1}^{n}(y_i - \hat{y}_i)^2}$. R-squared (R2) Score measures the proportion of variance in magnitude explained by the model, calculated as $R^2 = 1 - \frac{\sum_{i=1}^{n}(y_i - \hat{y}_i)^2}{\sum_{i=1}^{n}(y_i - \bar{y})^2}$.

\subsubsection{Results}

Table \ref{tab:regression_results} presents the performance of Random Forest Regressor on the magnitude prediction task.

\begin{table}[H]
\centering
\caption{Magnitude regression results}
\label{tab:regression_results}
\begin{tabular}{|l|c|c|c|}
\hline
\textbf{Model} & \textbf{RMSE} & \textbf{MAE} & \textbf{R2 Score} \\
\hline
Random Forest Regressor & 0.847 & 0.612 & 0.723 \\
\hline
\end{tabular}
\end{table}

The Random Forest Regressor achieved an RMSE of 0.847 magnitude units and R2 Score of 0.723, indicating that the model explains approximately 72.3\% of the variance in earthquake magnitude. The Mean Absolute Error of 0.612 suggests that on average, predictions deviate by about 0.6 magnitude units from actual values.

Analysis of prediction errors revealed that the model performs best for moderate magnitudes (M 2.0-4.0) where training data is most abundant. Performance degrades for extreme values, particularly very small magnitudes (M $<$ 0) and large magnitudes (M $>$ 6.0), due to limited training examples in these ranges.

Feature importance analysis confirmed that depth is the most predictive feature for magnitude, followed by the number of recording stations (nst) and spatial coordinates. This aligns with seismological understanding that larger earthquakes tend to occur at greater depths and are detected by more stations.

\subsection{Task 2: Depth Category Classification}

\subsubsection{Task Definition}

The objective of this task is to classify earthquakes into three depth categories following standard seismological conventions: Shallow earthquakes with depth less than 70 km, Intermediate earthquakes with depth between 70 and 300 km, and Deep earthquakes with depth greater than 300 km. Accurate depth classification is important for hazard assessment, as shallow earthquakes cause the most surface damage.

\subsubsection{Input Features}

The input features include latitude and longitude to capture regional depth patterns associated with different tectonic settings, number of recording stations (nst) which may correlate with earthquake characteristics, and azimuthal gap reflecting data quality. The actual depth value is excluded as the target variable.

\subsubsection{Model Implementation}

Two classification algorithms were implemented for this task. K-Nearest Neighbors (KNN) classifier leverages the spatial clustering of earthquakes by depth regime, where nearby earthquakes in feature space tend to have similar depths. The algorithm was configured with k=5 neighbors, distance-weighted voting, and Euclidean distance metric.

Decision Tree Classifier provides interpretable rules for depth classification based on geographic and quality features. The model was configured with maximum depth of 15, minimum samples per split of 20, and Gini impurity as the splitting criterion. Class weights were applied inversely proportional to class frequencies to handle imbalance.

\subsubsection{Evaluation Metrics}

Classification performance was evaluated using accuracy as the proportion of correct predictions, precision as the proportion of positive predictions that are correct, recall as the proportion of actual positives correctly identified (particularly important for the Shallow class due to hazard significance), and F1 Score as the harmonic mean of precision and recall. Macro-averaged metrics treat all classes equally regardless of size.

\subsubsection{Results}

Table \ref{tab:depth_classification_results} presents the performance of both classifiers on the depth category classification task.

\begin{table}[H]
\centering
\caption{Depth category classification results}
\label{tab:depth_classification_results}
\begin{tabular}{|l|c|c|c|c|}
\hline
\textbf{Model} & \textbf{Accuracy} & \textbf{Macro Precision} & \textbf{Macro Recall} & \textbf{Macro F1} \\
\hline
KNN (k=5) & 0.891 & 0.756 & 0.698 & 0.721 \\
Decision Tree & 0.903 & 0.782 & 0.724 & 0.748 \\
\hline
\end{tabular}
\end{table}

Decision Tree achieved slightly better performance than KNN with 90.3\% accuracy and 0.748 macro F1 score. Table \ref{tab:depth_perclass} shows per-class performance for the Decision Tree model.

\begin{table}[H]
\centering
\caption{Per-class performance for depth classification (Decision Tree)}
\label{tab:depth_perclass}
\begin{tabular}{|l|c|c|c|c|}
\hline
\textbf{Category} & \textbf{Support} & \textbf{Precision} & \textbf{Recall} & \textbf{F1 Score} \\
\hline
Shallow ($<$70 km) & 389,421 & 0.924 & 0.967 & 0.945 \\
Intermediate (70-300 km) & 38,567 & 0.712 & 0.623 & 0.665 \\
Deep ($>$300 km) & 5,276 & 0.708 & 0.582 & 0.639 \\
\hline
\end{tabular}
\end{table}

The Shallow class achieved excellent performance with 0.967 recall, which is critical for hazard assessment applications. Performance for Intermediate and Deep categories is lower due to class imbalance and the challenge of predicting depth from surface-observable features alone.

The Decision Tree model revealed interpretable patterns: earthquakes in certain latitude-longitude regions (corresponding to subduction zones) are more likely to be intermediate or deep, while events in continental regions are predominantly shallow. The number of recording stations also contributes to classification, as deeper earthquakes are often detected by more distant stations.

\subsection{Task 3: Seismic Hotspot Clustering}

\subsubsection{Task Definition}

The objective of this task is to identify natural clusters of seismic activity corresponding to tectonic features such as fault lines, subduction zones, and volcanic regions. Unlike supervised classification, clustering allows the data to reveal inherent spatial patterns without imposing prior assumptions about geographic boundaries.

\subsubsection{Input Features}

The input features include latitude and longitude for spatial clustering, and depth to distinguish between shallow crustal seismicity and deep subduction zone earthquakes. This three-dimensional feature space captures both horizontal and vertical distribution of seismic activity.

\subsubsection{Model Implementation}

Two complementary clustering algorithms were implemented. K-Means Clustering partitions earthquakes into a predefined number of clusters by minimizing within-cluster variance. The algorithm iteratively assigns points to nearest centroids and updates centroid positions. The optimal number of clusters was determined using the elbow method, which identifies the point where adding more clusters provides diminishing returns in variance reduction.

DBSCAN (Density-Based Spatial Clustering of Applications with Noise) identifies clusters based on point density without requiring a predefined number of clusters. Points in dense regions are grouped together, while points in sparse regions are classified as noise. DBSCAN naturally handles clusters of arbitrary shape, making it well-suited for detecting elongated fault line patterns. The algorithm was configured with epsilon (neighborhood radius) of 2.0 degrees and minimum samples of 100 points.

\subsubsection{Evaluation Metrics}

Clustering quality was evaluated using Silhouette Score, which measures how similar points are to their own cluster compared to other clusters. Scores range from -1 to 1, with higher values indicating better-defined clusters. The formula is $s(i) = \frac{b(i) - a(i)}{\max(a(i), b(i))}$, where $a(i)$ is the average distance to points in the same cluster and $b(i)$ is the average distance to points in the nearest other cluster.

Visual inspection of cluster boundaries against known tectonic features provides qualitative validation. Geographic interpretation by domain knowledge confirms whether identified clusters correspond to meaningful seismological structures.

\subsubsection{Results}

Table \ref{tab:clustering_results} presents the clustering performance metrics.

\begin{table}[H]
\centering
\caption{Clustering results}
\label{tab:clustering_results}
\begin{tabular}{|l|c|c|c|}
\hline
\textbf{Algorithm} & \textbf{Clusters} & \textbf{Silhouette Score} & \textbf{Noise Points} \\
\hline
K-Means (k=15) & 15 & 0.412 & 0 \\
DBSCAN & 23 & 0.387 & 12,847 \\
\hline
\end{tabular}
\end{table}

K-Means with 15 clusters achieved a Silhouette Score of 0.412, indicating moderately well-defined clusters. The elbow analysis suggested that 12-18 clusters provide a good balance between cluster compactness and interpretability.

DBSCAN identified 23 natural clusters with a Silhouette Score of 0.387. The algorithm classified 12,847 points (approximately 0.4\%) as noise, representing isolated earthquakes that do not belong to major seismic zones. The lower Silhouette Score compared to K-Means reflects the irregular shapes of clusters detected by DBSCAN.

Geographic analysis of the clustering results revealed meaningful patterns. K-Means clusters correspond to major seismic regions including the Pacific Ring of Fire segments (separate clusters for Japan-Kuril, Philippines-Indonesia, South America western coast, Central America, and Alaska-Aleutians), the Alpine-Himalayan belt (Mediterranean, Middle East, Himalayan regions), mid-ocean ridges, and continental rift zones.

DBSCAN clusters more closely follow linear fault traces and subduction zone geometries. The algorithm successfully identified the San Andreas Fault system, the Japan Trench, the Sumatra-Andaman subduction zone, and the Chile-Peru Trench as distinct elongated clusters.

\subsection{Summary of Results}

Table \ref{tab:summary_results} summarizes the key results across all three modeling tasks.

\begin{table}[H]
\centering
\caption{Summary of experimental results}
\label{tab:summary_results}
\begin{tabular}{|l|l|l|l|}
\hline
\textbf{Task} & \textbf{Best Model} & \textbf{Primary Metric} & \textbf{Result} \\
\hline
Magnitude Regression & Random Forest & R2 Score & 0.723 \\
Depth Classification & Decision Tree & Accuracy & 90.3\% \\
Hotspot Clustering & K-Means (k=15) & Silhouette Score & 0.412 \\
\hline
\end{tabular}
\end{table}

\subsection{Discussion}

The experimental results demonstrate that machine learning methods can effectively extract patterns from earthquake catalog data across different task types. The magnitude regression task achieved reasonable predictive performance, explaining over 70\% of magnitude variance using only location and quality features. This suggests that earthquake magnitude is partially predictable from observable characteristics, though significant unexplained variance remains.

The depth classification task achieved high accuracy, particularly for the critical shallow earthquake category. The interpretable Decision Tree model revealed geographic patterns consistent with tectonic understanding, where certain regions are associated with specific depth regimes.

The clustering analysis successfully identified major seismic zones without prior geographic knowledge. Both K-Means and DBSCAN revealed meaningful spatial patterns, with DBSCAN providing better representation of linear fault structures. These clusters can inform regional hazard assessment and seismic network planning.

Several limitations should be acknowledged. The regression and classification models are limited by the features available in catalog data, excluding waveform characteristics that could improve predictions. Class imbalance affects performance for minority categories in both classification and clustering. Model performance may vary across different geographic regions and time periods.
